\chapter{Shared mechanisms and genuine distinctions}\label{ch:synthesis}
\plain{The programs share some formulas, but they are not one interchangeable
experiment. Three pairwise-interacting points, eight centroid-coupled centers
and eight centers with stochastic response signs have different states,
interaction calculations and recording systems. Their evidence must be
compared at the same level of certainty.}

\section{A precise mechanism comparison}
Table~\ref{tab:mechanisms} summarizes the source-supported distinctions.
The complete side-by-side source locators are also supplied in
\filepath{CROSS_STUDY_MECHANISM_MAP.md}. Similar helper names do not
justify replacing one interaction law with another. \claim{C80}

\begin{table}[htbp]\centering\small
\caption{Mechanism map. State dimensions count real position/velocity
entries; the binary modes and external RNG/frame inputs are stated
separately. Source support: S01 lines 5--74; S02 lines 8--119;
S03 lines 10--167.}\label{tab:mechanisms}
\begin{tabularx}{\linewidth}{@{}l Y Y Y@{}}\toprule
Property & S01 & S02 & S03\\\midrule
State &12 reals&48 reals&48 reals + 8 bits\\
Positions &3 in 2D&8 in 3D&8 in 3D\\
Initial velocity &Zero&Components in $[-0.1,0.1)$&Same distribution as S02\\
Interaction &Two pairwise directions&Arithmetic centroid&Signed arithmetic-centroid response\\
Regularizer &Norm +0.001&Norm +0.01&Norm +0.01\\
Order &0,1,2 in place&0 through 7 in place&0 through 7 in place\\
Step multiplier &0.05&0.1&0.1\\
Coordinate period &2&10&10\\
H argument &Loop index $\times0.05$&Frame label $\times0.1$&Frame label $\times0.1$\\
Speed feedback &Absent&Shared GRF&Same GRF, also toggle gate\\
Randomness after setup &Absent&Absent&Gated per-body draws\\
Stored history &Post-sweep position log&No numeric history export&CSV on every callback\\
Display &Static connected trails&Moving tetrahedra&Moving tetrahedra with mode palettes\\\bottomrule
\end{tabularx}
\end{table}

S01's sum of regularized pairwise directions and S02's direction toward
an arithmetic mean are different nonlinear constructions. Averaging
positions before normalizing is not the operation written in S01.
The change of dimension, boundary period, regularizer and initial
velocity distribution further prevents interpreting their pictures as
a controlled comparison of one isolated mechanism.
\src{S01}{45--56}; \src{S02}{85--96}. \claim{C81}

There is a narrower exact relationship between S02 and S03:
given identical incoming numerical arrays and an identical frame label,
if every S03 post-gate mode is zero, its position/velocity update is
the same mathematical update as S02's. Both share the same mixed-state
centroid, H, GRF, norm regularizer, RPCO and wrap. This conditional
identity follows by substituting $b_i^+=0$ into
Equation~\eqref{eq:kick03}. It does not establish equal historical
initializers, random streams, callback schedules or trajectories.
\src{S02}{74--104}; \src{S03}{100--149};
\cite{study03}, SIMULATION\_02\_COMPARISON. \claim{C82}

\section{Sequential order is part of each model}
All three programs overwrite positions before later bodies are
processed. Consequently the mathematical state seen by a later
interaction is a mixture of new and old rows. Equations
\eqref{eq:mixed01} and \eqref{eq:centroid} preserve that dependence.
A simultaneous proposal/commit implementation would generally
calculate different directions; it would be a modified model,
not an innocent translation of the historical source.
\claim{C83}

The pictures show end-of-sweep states, so they hide this internal order.
S01's accepted replay supplies all within-sweep arrays and demonstrates
the dependency directly in events 1--3 and 505--507. S03's CSV exports
only complete callbacks; its mixed-state values used in accepted
residual checks are reconstructed from observed adjacent blocks.
Those different observation capabilities must remain visible in the
method, even when the algebraic notation is similar. \claim{C84}

\section{Persistent velocity and componentwise damping}
For any scalar component $u_a$, Equation~\eqref{eq:RPCO} gives
\begin{equation}\label{eq:damping}
\operatorname{sgn}\R(u)_a=\operatorname{sgn}u_a,\qquad
|\R(u)_a|\leq |u_a|,
\end{equation}
with strict magnitude reduction when $u_a\ne0$. This follows because
$1+\log(1+|u_a|)>1$ for a nonzero component. Different component
magnitudes receive different denominators, so the vector direction
can change. The magnitude comparison is with the \emph{candidate}
velocity after adding the drive, not necessarily with the previous
stored velocity. \src{S01}{29--31,54--56};
\src{S02}{44--46,91--96}; \src{S03}{47--49,123--126}.
\claim{C85}

Thus an inward response does not imply inward movement during the
same update. S01's near-contact and S03's delayed switch response
are two distinct observed examples. A zero instantaneous interaction
also does not erase a carried nonzero velocity. No source implements
a waiting state or a threshold that freezes every small velocity.
These source facts are not a theorem against approximate waiting
or approximate collective following over a particular finite window.
\claim{C86}

\section{A prescribed gain is not a route memory}
For nonnegative $t$, the shared H formula satisfies the elementary
bound
\begin{equation}\label{eq:Hbound}
|H(t)|\leq 2e^{-0.5t}+e^{-0.3t}+1.5e^{-0.7t}+e^{-0.4t}
\longrightarrow0\quad (t\to\infty).
\end{equation}
This is the triangle inequality applied to
Equation~\eqref{eq:H}, using $|\sin t|,|\cos t|\leq1$.
It does not say H is positive or that its magnitude decreases at
every sample. The accepted S02 analysis evaluates
$H(3\pi/2)\simeq-0.03600350414$ as a formula example; that time is
not claimed as a nominal frame label.
\cite{construction02}, TIME\_AND\_MOVEMENT. \claim{C87}

The bound concerns increasing arguments. It does not make a
restarting frame clock decay forever as callback count increases.
S03's original records demonstrate why that distinction matters:
at each return to label 0 the prescribed gain returns to 4.5.
Neither H nor GRF reads a saved itinerary. GRF is state feedback
through current mean speed; carried velocities are another form
of state persistence. A stored route that explicitly guides future
movement is absent from these three sources. \claim{C88}

\plain{There are several meanings of “memory” here. The current velocity
retains effects of the past. A log records the past for an observer.
H is a preset function of a supplied clock. These roles should be
named separately before relating them to a broader SRG time theory.}

\section{Coordinate geometry and measurement}
The coordinate rules identify opposite edges but do not add a
material container. Minimal wrapped displacements, native increments,
stored coordinate differences and screen distances are separate
quantities. S01 draws connectors between wrapped samples; S02 and S03
draw current shapes without trails. S03 adds a diagnostic whose
unsigned modulo differs from the displacement wrap used in the force.
\claim{C89}

The arithmetic centroid has its own coordinate-representative
dependence. A seam crossing can change the interaction target even
when the native movement is small. The measured S03 body 7 episode
is therefore not adequately described by calling every visible
effect a rendering artifact. Conversely, a large screen relocation
alone does not certify a large native movement. Equations
\eqref{eq:centroid}, \eqref{eq:seam} and \eqref{eq:baddetector}
make the three operations distinguishable. \claim{C90}

\section{What the experiments establish together}
The evidence supports three bounded conclusions. First, source order
and carried velocity explain why an interaction direction and a
movement direction can differ. Second, callback history and clock
labels are necessary for interpreting a time-dependent gain: a
200-frame movie is not automatically a 200-update record. Third,
the recording and rendering conventions can either reveal or hide
the distinction between near-contact, projected overlap, seam
crossing and a detector warning. \claim{C91}

It does not follow that the stochastic mode is responsible for
every visible difference between the two eight-body movies.
Their unknown initializers, distinct nominal durations, blitting
settings and encoding rates are uncontrolled differences.
Nor do the finite records establish asymptotic convergence,
a persistent orbit, a universal collective-following phenomenon
or a physical theory of gravity. The kinetic bookkeeping identity
describes the code's diagnostic; it is not a gravitational energy
equation. \claim{C92}
\FloatBarrier
